% P19.6: the Gick-Holyoak experiment series and Gentner's distinct theory.
\chapter{Analogy and Problem Transfer}
\label{sec:p19-analogy-and-problem-transfer}

An analogy can suggest a solution which a person would otherwise
miss. Several achievements are involved: representing a previous
case, noticing that it matters, finding useful correspondences,
and adapting its solution to the new constraints. Remembering
the previous case does not by itself establish that these other
steps occurred.

This chapter presents all five experiments in Gick and Holyoak's
1980 study and its auxiliary judgment tasks
\cite{GickHolyoak1980AnalogicalProblemSolving}. It then distinguishes
that paper's process account from Gentner's 1983 structure-mapping
theory \cite{Gentner1983StructureMapping}. Their proposed
representations are theoretical descriptions. A participant's
reported solution, a retrospective explanation and an analyst's
representation of the task provide different kinds of evidence.
The records follow \cref{def:p19-evidence-record}.

\section{The target problem, source stories and response codes}
\label{sec:p19-analogy-tasks}

\begin{hbdefinition}{def:p19-analogy-study-tasks}{The stipulated radiation task and its military analogues}
The target is a fictional problem adapted from Duncker: an
inoperable internal tumor must be destroyed while surrounding
healthy tissue is preserved. High-intensity rays destroy the
tumor but also damage tissue on their path; low-intensity rays
do not damage that tissue but are insufficient against the
tumor. The intended dispersion proposal sends low-intensity
rays from several directions so that they converge simultaneously
at the target. These are the task's stipulated constraints,
not measurements of a clinical treatment.

The principal source story concerns a fortress reached by
radial roads. A force large enough for the objective cannot
travel safely down one mined road, while small groups can.
The attack-dispersion ending divides the force among roads
and coordinates its arrival. The open-supply-route ending
uses one existing safe road. The tunnel ending constructs
a passage beneath the mined route. The first experiment
uses a common problem paragraph with these three endings.

The second experiment also uses a parade-dispersion story.
It retains the central fortress, radial roads and divided,
simultaneously arriving groups, but changes the goal to
an impressive parade. Loss of the general's rank replaces
the attack story's destructive consequences. The solution's
similar wording therefore occurs under different goals
and constraints.

The first three experiments explicitly suggest using the
source to help solve the target and permit looking back
at it. Experiments IV and V instead place the story in
a recall task and distinguish later instructions with
and without an explicit analogy hint.
\index{analogy!problem transfer}\index{analogy!noticing relevance}
\end{hbdefinition}

The original stimuli have three attack-story versions.
Version 1 is used in Experiment I. Version 2 is used
in II and, without its supplied ending, in III.
Version 3, entitled \emph{The General} for participants,
is used in IV and V. The original appendices preserve
their different wording. In particular the first two
versions connect village destruction to retaliation,
whereas the third describes destruction following the
mines' explosion. An abstract task summary should not
be mistaken for the literal wording of every version
\cite[pp.~351--354, appendices~I--IV]{GickHolyoak1980AnalogicalProblemSolving}.

\begin{hbdefinition}{def:p19-analogy-response-codes}{A proposed solution and a complete dispersion response}
A proposed solution is a suggestion recorded during
the task, including one that its author subsequently
rejects. From Experiment II onward, a complete
dispersion response contains all three of:
\begin{enumerate}
\item rays approaching from different directions;
\item low intensity on the separate paths;
\item simultaneous convergence.
\end{enumerate}
A partial dispersion response must include the first
feature but can omit the second or third, or allow
some damage to healthy tissue.
An open-passage response uses an existing route.
An operation response opens or inserts a route through
tissue. An accumulation response, used in Experiment
III, distributes ray applications over time.

Producing these categories does not establish endorsement
or successful satisfaction of the task constraints.
An operation conflicts with the intended prohibition
on operating; the source also records that some people
read that prohibition as excluding tumor removal
specifically. That interpretation difference matters
when judging a response.
\end{hbdefinition}

For the reported analyses below, \(G^2(d)\) denotes a
likelihood-ratio contingency-table statistic with \(d\)
degrees of freedom. \(F(d_1,d_2)\) denotes a variance-ratio
statistic with the displayed numerator and denominator
degrees of freedom; \(t\) denotes the reported mean-comparison
statistic. A \(p\)-value refers to the source's stated
analysis, with the interpretation principles of
\cref{sec:p19-evidence-methods}.

\section{Experiment I: prompted transfer and verbal protocols}
\label{sec:p19-analogy-experiment-one}

Forty introductory-psychology undergraduates at the University
of Michigan participated for course credit, ten in each
of four conditions: attack-dispersion, open supply route,
tunnel, and no prior story. Assignment is reported;
random assignment is not specified for this experiment.
Participants were tested individually and recorded on tape.
A candle problem first familiarized them with thinking aloud.
Story participants read the story aloud, summarized its gist,
and were instructed to try using it as a hint, while being
told that it was not necessary to do so.

The experimenter could request clarification, including a
drawing, and could ask unsuccessful participants to reread
the instructions and then to use the story. Two independent
scorers coded the transcripts; condition blinding is not
reported here. Multiple proposal categories can occur for
one person. The original Table 3 reports:
\[
\begin{array}{l|r|rrr}
 \text{prior story}&N&\text{dispersion}&\text{open passage}&\text{operation}\\ \hline
 \text{attack-dispersion}&10&100\%&10\%&30\%\\
 \text{open supply route}&10&10\%&70\%&50\%\\
 \text{tunnel}&10&20\%&30\%&80\%\\
 \text{none}&10&0\%&20\%&50\%
\end{array}
\]
One of the seven open-passage responses in the open-route
condition proposed an unspecified passage without identifying
a concrete route. Operation proposals include opening a path
or inserting a tube through tissue
\cite[pp.~319--321]{GickHolyoak1980AnalogicalProblemSolving}.

The reported dispersion comparison is \(G^2(3)=33.9\),
\(p<.001\); attack-dispersion versus the other groups gives
\(G^2(1)=30.9\), \(p<.001\). The open-route condition's
70\% open-passage rate versus 20\% across the other
conditions gives \(G^2(1)=8.21\), \(p<.01\).
The tunnel condition's 80\% operation rate versus about
43\% across the other conditions gives \(G^2(1)=4.29\),
\(p<.05\). The source uses the likelihood-ratio
contingency-table statistic \(G^2\), rather than treating
every proposed idea as an independent participant.

Of the ten attack-dispersion participants, seven produced
dispersion without a further prompt and three after a
further prompt. All ten had already received the initial
instruction to use the analogy. The 100\% result is
therefore prompted transfer, and the seven early responses
are not a no-hint experimental condition.

The protocols offer process examples with explicit limits.
Participant S15 developed an initially sequential
low-intensity proposal into multidirectional convergence
after a reminder. S19 mapped an open route abstractly
without producing a concrete passage. S24 recognized
the tunnel parallel but encountered the operation
constraint. S28 reported a previously encountered
lightbulb-and-laser example, while also demonstrating
knowledge of the military analogy; the source does
not establish which example caused the solution.
The published excerpts are selected illustrations,
not an exhaustive classification of all participants'
mental processes
\cite[pp.~325--331, tables~5--8]{GickHolyoak1980AnalogicalProblemSolving}.

Four of seven open-passage solvers in the open-route group
and four of eight operation proposers in the tunnel group
received additional prompts before those proposals.
Controls received more encouragement to talk but fewer
requests to clarify solutions. The authors identify
experimenter interaction as a reason to use a noninteractive
procedure in Experiment II.

\subsection{Other proposals and judgments of one's own solutions}

After excluding each story group's principal analogous
category, mean numbers of other proposals were 1.10,
1.60 and 2.70 for attack, open route and tunnel,
respectively, versus 2.00 for controls. Controls had
no analogous category removed. Non-ray proposals,
averaging 0.4 per person, were excluded in all groups.
The reported analysis is \(F(3,36)=4.17\), \(p<.025\);
the Newman--Keuls comparison distinguished only tunnel
from attack, \(p<.01\).

Proposals to reduce healthy tissue's sensitivity or
shield it occurred for 30\% of controls, 10\% of the
tunnel group and none in the other two groups.
The paper explicitly treats these small counts as
insufficient for a reliable statistical conclusion.
They cannot establish that analogy invariably suppresses
alternative ideas.

Participants also rated their own proposals for
practicality and creativity on two seven-point scales,
where 1 was the maximum. Among the seven attack
participants who had both dispersion and another
proposal, mean practicality ratings were 1.43 versus
4.64, \(F(1,6)=27.2\), \(p<.01\). Creativity ratings
were 3.43 versus 3.64, \(F<1\). Parallel within-person
comparisons in the other story groups were not
significant. Pooled practicality means were 2.0 for
dispersion, 3.9 for open passage and 5.4 for operations.
These are judgments by the people who produced the
respective proposals; they are not randomized comparisons
of treatment effectiveness
\cite[pp.~323--325]{GickHolyoak1980AnalogicalProblemSolving}.

\section{Auxiliary judgments of source prompts and practicality}
\label{sec:p19-analogy-auxiliary}

The source additionally reports a story-rating task
with 17 undergraduates per story. The three groups
in Table 4 comprise 51 participants unfamiliar with
a solution to the target. They read a story, the
target, and six proposed solutions in random order.
They judged how likely the story was to suggest
each proposal on a nine-point scale, with 9 the
maximum, and were told to disregard practicality.
This scale runs in the opposite direction from
Experiment I's seven-point self-ratings.

The six proposal codes below denote:
\(1\), simultaneous dispersion;
\(2\), an existing esophageal passage;
\(3\), a tube inserted through tissue;
\(4\), an incision exposing the tumor;
\(5\), repeated low-intensity applications; and
\(6\), intermediate-intensity rays.
The reported mean ratings are:
\[
\begin{array}{l|rrrrrr}
 \text{story}&1&2&3&4&5&6\\ \hline
 \text{attack-dispersion}&8.41&2.53&3.71&2.94&4.88&3.47\\
 \text{open supply route}&3.59&6.65&5.82&4.58&2.94&2.29\\
 \text{tunnel}&5.12&4.47&6.47&5.35&2.71&1.82
\end{array}
\]
The story-by-proposal comparison gives
\(F(10,240)=8.84\), \(p<.001\).
For attack, dispersion was rated above every other
proposal, \(p<.01\). For the open-route story, proposals
2--4 did not differ significantly from each other
but exceeded the other proposals, \(p<.05\).
For tunnel, the two operation proposals, dispersion
and open passage did not differ significantly from
each other; the operation proposals exceeded 5 and 6,
\(p<.05\). These comparisons used Newman--Keuls tests.

The parade story was an additional 17-person condition
of this same rating task. Its comparison with the
17-person attack condition uses 34 participants,
including the attack group already represented in
Table 4; it is not another independent sample of
34 people. Across the six proposals the source reports
a main effect \(F<1\) and interaction
\(F(5,160)=1.32\), \(p>.25\). It does not print the
six parade means. The absence of a detected rating
difference does not establish equivalence of the
stories' effects on generating solutions
\cite[pp.~321--323, 333]{GickHolyoak1980AnalogicalProblemSolving}.

A separate group of twenty undergraduates unfamiliar
with the problem chose which of dispersion and an
esophageal route seemed more practical. Fifteen chose
dispersion, with reported sign-test \(p<.05\);
eleven of those fifteen mentioned possible injury
to the intervening tissue. This is a forced-choice
judgment task, distinct from producing a solution
and from evaluating actual medical outcomes.

\section{Experiment II: analogous and disanalogous problem statements}
\label{sec:p19-analogy-experiment-two}

Experiment II used booklets in five introductory
psychology classes. There were initially 143 participants:
47 read attack-dispersion version 2, 46 read parade-dispersion,
and 50 received no story. One parade participant with
prior knowledge of the target and dispersion solution
was excluded, leaving 142 for analysis.

The attack and parade stories were matched in length
and wording where possible. Their solution paragraphs
were the same except for the final goal statement.
Their problem constraints and the role played by
convergence differed. Participants summarized the
source in writing, could look back at it, and
received the initial suggestion that it might help.
They wrote proposals in order, including later-rejected
ideas. No interactive clarification prompts were used.
Two independent scorers were blind to condition;
disagreements were resolved by discussion.

The complete-versus-partial coding of
\cref{def:p19-analogy-response-codes} was introduced
to handle incomplete written descriptions.
The reported Table 10 percentages are:
\[
\begin{array}{l|r|rrr}
 \text{condition}&N&\text{complete}&\text{partial}&\text{reported total}\\ \hline
 \text{attack}&47&57\%&19\%&76\%\\
 \text{parade}&45&31\%&18\%&49\%\\
 \text{control}&50&8\%&0\%&8\%
\end{array}
\]
These are the source's rounded displays; the total
column adds the displayed component percentages.
It should not be treated as an exact fraction
with denominator \(N\). All four control solvers
gave complete responses.

The reported overall any-dispersion comparison is
\(G^2(2)=53.1\), \(p<.001\). Attack versus parade
has \(G^2=7.70\), \(p<.01\), and the combined story
conditions versus control have \(G^2(1)=45.5\),
\(p<.001\). For the stated two-condition comparison
of any dispersion against none, the contingency
table has one degree of freedom. Distinguishing
complete, partial and absent responses would instead
define a different comparison
\cite[pp.~332--335]{GickHolyoak1980AnalogicalProblemSolving}.

Dispersion tended to occur before the participant's
median proposal position, with reported sign-test
\(p<.025\). Mean counts of other proposals were
1.17 for attack, 1.40 for parade and 2.12 for controls,
\(F(2,139)=12.7\), \(p<.001\). The reported
Newman--Keuls comparison distinguished controls
from the combined story groups, \(p<.01\).
Tissue-protection proposals occurred for 2\%,
16\% and 28\%, respectively,
\(G^2(2)=14.6\), \(p<.01\); parade exceeded
attack in that category, \(G^2(1)=5.78\),
\(p<.02\).

\subsection{Retrospective helpfulness}

The following Table 11 percentages are conditional
on both story and produced response. Within each
triple the order is not helpful, somewhat helpful,
very helpful:
\[
\begin{array}{l|rrr|rrr|rrr}
 &\multicolumn{3}{c|}{\text{complete}}
 &\multicolumn{3}{c|}{\text{partial}}
 &\multicolumn{3}{c}{\text{none}}\\
 \text{attack}&0&22&78&11&22&67&64&27&9\\
 \text{parade}&0&29&71&0&38&63&57&39&4
\end{array}
\]
Rounding can make a triple total 101. These percentages
are not percentages of all recruited participants.
The source reports no reliable attack/parade difference
in the rating patterns and explicitly raises hindsight
as an alternative to accurate memory of the solution
process. Retrospective helpfulness is therefore not a
direct measurement of which correspondence generated
the solution.

\section{Experiment III: solutions generated by the participants}
\label{sec:p19-analogy-experiment-three}

Forty-six students in two introductory classes received
booklets. One with prior knowledge of the target's
dispersion solution was excluded, leaving 45.
They first received only the problem paragraph of
attack-dispersion version 2 and generated their own
solutions. They then attempted the radiation problem,
with an explicit suggestion to use those solutions
and permission to look back. Two independent scorers
coded the responses and resolved disagreements by
discussion; condition blinding is not specified.

Twenty-two of 45 participants generated the gist of
military dispersion. Nine of those 22 subsequently
gave target dispersion, five complete and four partial.
Of the other 23 participants, three gave target
dispersion, two complete and one partial:
\[
\begin{array}{l|r|rr|r}
 \text{self-generated military dispersion}&N&
    \text{complete target}&\text{partial target}&\text{any target}\\ \hline
 \text{yes}&22&5&4&9\\
 \text{no}&23&2&1&3
\end{array}
\]
The reported association is \(G^2(1)=4.61\),
\(p<.05\). The successful-source subgroup's
41\% target rate exceeds the earlier control
rate of 8\%, \(G^2(1)=10.4\), \(p<.01\),
and is lower than the earlier supplied-attack-story
rate of 76\%, \(G^2(1)=8.25\), \(p<.01\)
\cite[pp.~337--340]{GickHolyoak1980AnalogicalProblemSolving}.

Membership in the two rows was not experimentally
assigned: participants themselves generated or failed
to generate the source solution. The association can
therefore reflect prior ability or strategy as well
as transfer. The comparisons with Experiment II
also concern different samples and procedures.
They do not isolate the causal effect of generating
rather than reading a solution.

Participants gave a mean of 3.04 military proposals.
Ten of the 22 dispersion proposals omitted clear
simultaneity or added regrouping before attack;
four of those ten participants transferred.
The source found no detectable effect of that
quality distinction. It does not establish that
every incomplete source solution is equally effective.
The helpfulness questionnaire is described as mostly
positive among those solving both problems, with
a missing response; its wording does not provide an
unambiguous numerical allocation for that nine-person
subgroup.

Among the remaining participants the source reports
58\% rating the source not helpful. Two of the
three target-dispersion solvers who had not produced
military dispersion rated it somewhat or very helpful;
the third did not answer.

Another source proposal sent small groups successively
down one route. Of 24 participants giving it, ten
proposed temporally accumulated radiation, versus
three of the other 21, \(G^2(1)=4.28\), \(p<.05\).
Those rates are about 42\% and 14\%.
The military proposal categories can overlap, as can
the target categories. The suggestions were recorded
even if later rejected, so these counts do not show
that participants believed the accumulated treatment
would work.

\section{Experiment IV: remembering a source and noticing its relevance}
\label{sec:p19-analogy-experiment-four}

Twenty-seven paid undergraduates from the Human
Performance Center subject pool were randomly assigned
to hint (\(N=12\)) or no hint (\(N=15\)).
In small groups they memorized three stories in a
fixed order: \emph{The Wine Merchants}, \emph{The General}
and \emph{The Identical Twins}. Each story received
three minutes of study, was collected, and was
followed by up to fifteen minutes of recall.
All completed recall within that time. The middle
story was the relevant attack-dispersion version 3.
The other stories used a barrel raft and a substituted
twin, respectively, rather than the convergence plan.

After a three-to-five-minute break participants wrote
target solutions. The hint condition suggested that
one earlier story might help; that sentence was absent
in the no-hint condition. Two independent scorers
agreed completely on the dispersion coding.
Eleven of twelve hinted participants produced complete
dispersion. In the no-hint group three of fifteen
produced dispersion, one complete and two partial,
\(G^2(1)=15.5\), \(p<.001\):
\[
\begin{array}{l|r|rr|r}
 &N&\text{complete}&\text{partial}&\text{any dispersion}\\ \hline
 \text{hint}&12&11&0&11\\
 \text{no hint}&15&1&2&3
\end{array}
\]

Twelve no-hint participants reported that using the
stories had not occurred to them. The other three
reported using the attack story and rated it very
helpful, but only one produced dispersion. The two
partial solvers said that they had not considered
the stories. Thus three produced at least part of
the solution, three reported story use, and only one
belonged to both groups. These are different measures
\cite[pp.~340--343]{GickHolyoak1980AnalogicalProblemSolving}.

In the hint group ten rated the attack story very
helpful, one somewhat helpful and one did not answer.
Half rated the wine story not helpful; ten of eleven
respondents rated the twins story not helpful.
The critical story was divided into 43 gist units
for recall coding. Mean recall was 32.08 units in
the hint group and 33.53 in the no-hint group,
\(t<1\). Every participant recalled the country
setting and the two units describing radial roads.
The absence of a detected mean recall difference
does not establish statistical equivalence, but
the recorded recall and later prompted performance
show why recall and spontaneous relevance detection
must be distinguished in interpreting this task.

\section{Experiment V: a single story before or during problem solving}
\label{sec:p19-analogy-experiment-five}

Forty-seven paid undergraduates from the same subject
pool were assigned to story first (\(N=17\)), story
second (\(N=20\)) or incubation control (\(N=10\)).
Random assignment is not explicitly stated.
Story first involved studying and recalling only
the attack story, then an unhinted target attempt.
Participants were asked whether using the story had
occurred to them; those answering no were instructed
to use it for additional proposals.

Story second began with a ten-minute target attempt,
followed by attack-story study and recall, another
unhinted target attempt, and a final attempt after
a hint. The incubation control followed that sequence
with the unrelated twins story instead. The procedures
therefore distinguish initial solution, solution after
the intervening story, and solution after the later
instruction to use it.

Table 12 classifies the stage at which dispersion first
appeared. Entries in a row describe mutually exclusive
groups, not independent trials with a newly recruited
sample at every stage:
\[
\begin{array}{l|r|rrrr}
 \text{condition}&N&\text{before story}&
 \text{after, no hint}&\text{after, with hint}&\text{never}\\ \hline
 \text{story first}&17&\text{not applicable}&41\%&35\%&24\%\\
 \text{story second}&20&10\%&35\%&30\%&25\%\\
 \text{incubation control}&10&10\%&0\%&0\%&90\%
\end{array}
\]
The corresponding counts are \(7,6,4\) for story first;
\(2,7,6,5\) for story second; and \(1,0,0,9\) for
the control. One story-first solution was partial;
the other reported solutions were complete
\cite[pp.~343--346]{GickHolyoak1980AnalogicalProblemSolving}.

Six of the seven story-first participants who solved
without a hint reported using the story. On the
authors' self-report interpretation, successful
spontaneous use was about 35\% in each story condition.
The proportions who said using the story had occurred
to them were 47\% and 40\%, respectively.
The experiment supplied no detected advantage from
placing the story after an initial target attempt.
This finding is not an equivalence result about all
ways of encountering an analogy during problem solving.

Across the two relevant-story conditions, sixteen of
37 participants solved before the final hint and
28 of 37 eventually solved. The source reports these
as 43\% and 76\%. The sixteen include two who solved
\emph{before receiving the story}; fourteen first
solved after the story without the hint. Thus the
43\% figure cannot be read as the proportion whose
solution was demonstrably caused by spontaneously
using that story.

The reported 43\% versus Experiment IV's 20\% no-hint
comparison is \(G^2(1)=2.64\), \(p=.10\).
It also changes both the number of distractor stories
and the relevant story's serial position, so it does
not isolate a memory-load effect. Eventual production
in the two relevant-story conditions versus the
incubation control gives \(G^2(1)=15.0\), \(p<.001\).
The control's sole solver succeeded before the break;
no additional control participant solved afterward.
That small condition does not establish a universal
absence of incubation effects.

\section{The two theoretical presentations}
\label{sec:p19-analogy-theories}

Gick and Holyoak distinguish building representations,
mapping the source onto the target, and generating
the corresponding target solution. Their Figure 1
and Table 1 are approximate analyst representations
with propositional and higher-order links, including
inferences not literally present in the stories.
They allow correspondences between semantically
related predicates, such as capture and destroy,
and discuss different levels of abstraction.
A solution-focusing strategy might map just enough
of a source ending to generate a proposal; a more
complete mapping might also use the problem constraints.

Their own example is not a perfect isomorphism.
The army generally corresponds to the rays, but the
army occupies a threatened-object role in the detailed
mine representation while rays occupy a destructive-
instrument role in the target representation. The
radial roads also have no explicitly stated target
counterpart. Choosing a coarser representation can
hide one mismatch while omitting information needed
for adaptation. The experiments do not recover a
unique internal representation for every participant
\cite[pp.~308--316, 346--351]{GickHolyoak1980AnalogicalProblemSolving}.

Gentner's structure-mapping account distinguishes
object attributes, relations between objects, and
higher-order relations whose arguments include
propositions. Given object correspondences, it
generates candidate target predicates by substitution.
It gives priority to relational structure over
object attributes and favors connected systems
of relations over isolated matches. This last
preference is the systematicity principle
\cite[pp.~155--165]{Gentner1983StructureMapping}.

The syntactic distinctions are relative to the
chosen representation. Gentner explicitly notes
that a many-place relation can be re-encoded
as a one-place predicate, and that comparisons
with a prototype can function psychologically
as attributes. Her presentation assumes predicate
identity at a suitable component level, whereas
Gick and Holyoak explicitly discuss semantic
transformations between predicates. These are
distinct assumptions about representation.

The main comparison types in Gentner's account are:
\[
\begin{array}{l|cc}
 \text{comparison}&\text{object-attribute overlap}&
                   \text{relational overlap}\\ \hline
 \text{literal similarity}&\text{substantial}&\text{substantial}\\
 \text{analogy}&\text{limited}&\text{substantial}\\
 \text{mere appearance}&\text{substantial}&\text{limited}\\
 \text{anomaly}&\text{limited}&\text{limited}
\end{array}
\]
An abstraction also transfers a relational structure,
but its base objects are already general rather
than concrete objects whose individual attributes
must be left behind. The source treats these
contrasts as continua. Metaphors can mix relational
and attribute comparisons; not every metaphor
obeys one uniform analogy rule.

Systematicity is not a rule to copy every relation
or to prefer only causal information. In the source's
idealized central-force example, a mass relation
participates in the connected system while a
temperature relation does not. Its water/heat
example distinguishes a connected process description
from an isolated spatial description. These illustrate
the proposed selection principle, not empirical
validation of every projected physical statement.

The 1983 account also qualifies processing order:
object correspondences need to be settled before
the relational mapping is finished, not necessarily
before it begins. Relations can help determine
the object assignments. Its one-to-one explanatory
examples do not exclude every other kind of analogy.
Its analogical-shift sequence---literal comparison,
analogy, then an abstract law---is a conjecture,
illustrated by an informal observation of one
learner. The paper's empirical-support section
refers to other studies without supplying their
complete samples, coding details or numerical
results; it is not an additional fully reported
experiment in the present corpus
\cite[pp.~158, 164--170]{Gentner1983StructureMapping}.

\section{Explicit mapping and feasibility benchmarks}
\label{sec:p19-analogy-benchmarks}

The following constructions make particular comparison
rules precise. They are stipulated handbook models.
They are not assertions that the participants used
these symbols or calculations, and they do not define
a unique psychological algorithm.

\begin{hbdefinition}{def:p19-analogy-expression-map}{Objects, relation expressions and substitution}
Let \(B,T\) be finite object sets. A declared expression
language has object attributes of arity one, object
relations of specified finite arity, and higher-order
constructors with specified argument sorts.
Every constructor returns a proposition expression.
Syntax is free: no equations identify distinct expression
trees before any later interpretation.
Expressions are finite well-founded trees built from
objects by these constructors; an object and a
proposition expression have different sorts.

For example, with objects \(a,b,c\), object relations
\(R,S\), attribute \(A\), and a binary constructor
\(\mathsf{Supports}\) taking proposition expressions,
one can form
\[
 A(b),\qquad R(a,b),\qquad S(b,c),\qquad
             \mathsf{Supports}(R(a,b),S(b,c)).
\]
The name \(\mathsf{Supports}\) is a formal constructor
here, with no unprovided causal semantics.
An object correspondence \(m:B\to T\) induces a
candidate substitution \(\tau_m\) by replacing objects
with their images and leaving each constructor
and its ordered argument positions fixed.
\index{analogy!structural correspondence}\index{analogy!systematicity}
\end{hbdefinition}

\begin{hbtheorem}{thm:p19-analogy-structural-map}{proved=annex:p19.analogy-structural-map}{What a consistent object substitution preserves}
There is exactly one such recursive substitution
\(\tau_m\). It preserves sorts, constructor arities,
argument positions, nesting and every repeated
occurrence of a given object. If \(m\) is injective,
then \(\tau_m\) is injective on expressions.
With \(m(a)=x,m(b)=y,m(c)=z\), the displayed
connected expression becomes
\[
        \mathsf{Supports}(R(x,y),S(y,z)).
\]
One cannot map the occurrence of \(b\) in \(R(a,b)\)
to \(y\) while mapping its occurrence in \(S(b,c)\)
to a different \(y'\) under this single correspondence.
Discarding the independent attribute \(A(b)\)
does not disturb the displayed connected expression.
\end{hbtheorem}

\begin{hbdefinition}{def:p19-analogy-threshold-model}{A shared threshold model for dispersed resources}
Fix a route capacity \(h>0\), a target requirement
\(q>h\), and a positive integer \(k\) of available
independent routes. Amount \(x_j\) on route \(j\)
must satisfy \(0\le x_j\le h\). At the target,
simultaneously arriving amounts add, and success
requires \(\sum_{j=1}^k x_j\ge q\).
Independent routes here means that using one
does not add load to another before convergence.
It is a condition of this mathematical model,
not probabilistic independence.
\end{hbdefinition}

\begin{hbtheorem}{thm:p19-analogy-threshold-transfer}{proved=annex:p19.analogy-threshold-transfer}{The feasibility condition and the assumptions a plan transfers}
A feasible dispersed allocation exists exactly when
\(kh\ge q\). The minimum number of routes is
\(\lceil q/h\rceil\). Scaling \(h,q\) and every
allocation by the same positive factor preserves
feasibility. For \(h=4,q=10,k=3\), the equal
allocation \(x_j=10/3\) is feasible, while
putting all ten units on one route is not.

Copying that three-route plan to a target with
only two routes fails the requirement.
Even with three routes, if combined effect is
the maximum individual amount rather than
their sum, no safe allocation succeeds.
Thus central location and visual resemblance
alone do not establish the premises on which
the successful transfer depends.
\end{hbtheorem}

\begin{hbtheorem}{thm:p19-analogy-candidate-truth}{proved=annex:p19.analogy-candidate-truth}{A structurally available candidate need not follow logically}
Let the known target facts be \(R(x,y)\) and
\(S(y,z)\), and let \(P\) be a further binary
predicate with no supplied constraint connecting
it to those facts. The same object substitution
which maps source relations can generate the
candidate \(P(x,z)\) from \(P(a,c)\).
Nevertheless the known target facts entail
neither \(P(x,z)\) nor its negation.
There are target structures agreeing on all
the known facts and the object correspondence
but disagreeing on \(P(x,z)\).

Likewise, the formal presence of a higher-order
constructor in a mapped expression establishes
its syntax, not the truth of a causal,
explanatory or physical interpretation assigned
to that constructor.
\end{hbtheorem}

\section{What the evidence supports}
\label{sec:p19-analogy-evidence-scope}

The experiment series shows that a supplied remote
story can influence which target proposals these
participants produce. Its different designs separate
some effects of hints, source problem structure,
self-generated solutions and intervening recall.
The strength of a claim must follow the particular
design: Experiment IV explicitly reports random
assignment, Experiment III's successful-source
subgroups are self-selected, and comparisons across
experiments also change participants and procedures.

The series uses one target problem and a limited
set of stories. It does not establish transfer
rates for arbitrary scientific discoveries,
multi-step proofs or everyday decisions.
The source leaves open how spatial imagery,
different abstraction levels and search through
a large memory contribute. Selected protocols
and post-task helpfulness reports constrain an
interpretation but do not uniquely identify
an internal process.

The two theoretical accounts help articulate
candidate correspondences and missing assumptions.
The mathematical benchmarks show what a specified
mapping or constraint system guarantees.
Accepting a proposed real-world solution still
requires evidence that the target satisfies
the relevant constraints. Structural coherence
and observed transfer are each insufficient
to establish that conclusion on their own.
